\documentclass[12pt, a4paper]{article}

\usepackage[utf8]{inputenc}
\usepackage[T1]{fontenc}
\usepackage[brazil]{babel}
\usepackage{geometry}
\usepackage{graphicx}
\usepackage{booktabs}
\usepackage{caption}
\usepackage{subcaption}
\usepackage{hyperref}
\usepackage{amsmath}
\usepackage{float}
\usepackage{xcolor}
\usepackage{listings}
\usepackage{parskip}

\geometry{margin=2.5cm}

\hypersetup{
    colorlinks=true,
    linkcolor=black,
    urlcolor=blue,
    citecolor=black
}

\lstset{
    basicstyle=\ttfamily\small,
    backgroundcolor=\color{gray!10},
    frame=single,
    breaklines=true,
    keywordstyle=\color{blue},
    commentstyle=\color{green!60!black},
}

% ─────────────────────────────────────────────────────────────────────────────
\title{
    \textbf{Detecção de Árvores \textit{Cecropia} (Embaúba) em\\
    Mosaico Aéreo com Visão Computacional Clássica}\\[0.5em]
    \large INE5443 — Reconhecimento de Padrões \\ UFSC
}
\author{Augusto H. V. Guerner}
\date{Junho de 2026}
% ─────────────────────────────────────────────────────────────────────────────

\begin{document}

\maketitle
\tableofcontents
\newpage

% ─────────────────────────────────────────────────────────────────────────────
\section{Introdução}
% ─────────────────────────────────────────────────────────────────────────────

A \textit{Cecropia} (popularmente conhecida como Embaúba) é um gênero de árvore pioneira
amplamente distribuída em ambientes tropicais e subtropicais brasileiros. Sua
presença em áreas de vegetação secundária é um indicador de regeneração florestal,
tornando-a de interesse em estudos de cobertura vegetal e dinâmica de paisagem.

Este trabalho apresenta um pipeline de visão computacional clássica — sem uso de
aprendizado profundo — para detecção automática de indivíduos de \textit{Cecropia}
em um mosaico aéreo de alta resolução. A abordagem utiliza segmentação por espaço
de cor HSV combinada com operações morfológicas e filtragem geométrica por contorno,
processada em paralelo sobre 992 \textit{tiles} do mosaico.

% ─────────────────────────────────────────────────────────────────────────────
\section{Dataset}
% ─────────────────────────────────────────────────────────────────────────────

O mosaico aéreo foi dividido em tiles JPEG de resolução fixa. A
Tabela~\ref{tab:dataset} resume as características do dataset.

\begin{table}[H]
\centering
\caption{Características do dataset.}
\label{tab:dataset}
\begin{tabular}{ll}
\toprule
\textbf{Parâmetro} & \textbf{Valor} \\
\midrule
Total de tiles & 992 \\
Resolução por tile & 2048 $\times$ 2048 px \\
Sobreposição entre tiles & 512 px (25\%) \\
Sistema de referência (CRS) & EPSG:31982 (UTM Zona 22S) \\
Resolução espacial & $\approx$ 0,00848 u/px \\
Tiles pretos (fora da área mapeada) & 452 (ignorados automaticamente) \\
Tiles efetivamente processados & 590 \\
\bottomrule
\end{tabular}
\end{table}

Os tiles completamente pretos correspondem a regiões fora dos limites do mosaico
(padding do recorte) e são descartados antes do processamento com base na
intensidade média da imagem.

% ─────────────────────────────────────────────────────────────────────────────
\section{Pipeline de Detecção}
% ─────────────────────────────────────────────────────────────────────────────

O pipeline é composto por seis etapas sequenciais aplicadas a cada tile
individualmente. A Figura~\ref{fig:pipeline} ilustra todas as etapas sobre o
tile \texttt{tile\_0399.jpg}.

\begin{figure}[H]
    \centering
    \includegraphics[width=\textwidth]{../output/passo_a_passo/tile_0399.png}
    \caption{Visualização das 8 etapas do pipeline sobre \texttt{tile\_0399.jpg}
             (28 detecções). Da esquerda para a direita, de cima para baixo:
             imagem original, blur gaussiano, canal H, canal S, canal V,
             máscara bruta, máscara limpa, resultado final.}
    \label{fig:pipeline}
\end{figure}

\subsection{Etapa 1 — Suavização (Gaussian Blur)}

A imagem original em BGR é suavizada com um filtro gaussiano de kernel $9 \times 9$
e $\sigma = 0$ (calculado automaticamente a partir do tamanho do kernel). O objetivo
é reduzir ruído de alta frequência introduzido pela compressão JPEG e variações de
iluminação local, evitando que pequenas manchas de cor disparem falsas ativações
na etapa de segmentação.

\subsection{Etapa 2 — Segmentação HSV (\texttt{inRange})}

A imagem suavizada é convertida para o espaço de cor HSV (\textit{Hue, Saturation,
Value}). Esse espaço é mais robusto a variações de iluminação do que RGB, pois
separa a informação de matiz (cor) da luminosidade.

A segmentação é feita com um limiar fixo definido empiricamente para a tonalidade
verde-clara característica da folhagem de \textit{Cecropia}:

\begin{equation}
    \text{HSV}_{lower} = [44,\; 144,\; 104], \qquad
    \text{HSV}_{upper} = [58,\; 231,\; 203]
\end{equation}

Pixels cujos valores HSV estão dentro desse intervalo recebem valor 255 na máscara
binária (\textit{máscara bruta}); os demais recebem 0.

\subsection{Etapa 3 — Fechamento Morfológico (CLOSE)}

A máscara bruta contém fragmentos desconectados correspondentes a partes visíveis
da copa entre as folhas. A operação de fechamento morfológico com um elemento
estruturante elíptico de $25 \times 25$ pixels preenche essas lacunas, unindo
fragmentos próximos em regiões coesas que representam copas individuais.

\subsection{Etapa 4 — Abertura Morfológica (OPEN)}

Após o fechamento, pequenas regiões isoladas (ruído) que não correspondem a copas
reais são eliminadas pela abertura morfológica com elemento elíptico de $9 \times 9$
pixels. Essa etapa remove ativações espúrias de pequena área sem afetar as regiões
maiores já consolidadas pelo fechamento.

\subsection{Etapa 5 — Filtragem de Contornos por Área}

Os contornos externos da máscara limpa são extraídos com \texttt{cv2.findContours}.
Contornos com área inferior a $10.000\;\text{px}^2$ são descartados — threshold
que corresponde aproximadamente a uma copa de $\sim 0,7 \times 0,7\;\text{m}$ na
resolução do mosaico, eliminando detecções de vegetação rasteira e artefatos
residuais.

\subsection{Etapa 6 — Convex Hull}

Para cada contorno aprovado, calcula-se o fecho convexo (\textit{convex hull}).
Isso regulariza a forma da detecção, eliminando reentrâncias causadas por sombras
ou sobreposição de folhas, e aproxima melhor o perímetro circular da copa da
\textit{Cecropia} vista de cima. O hull é desenhado sobre a imagem original para
visualização.

% ─────────────────────────────────────────────────────────────────────────────
\section{Resultados}
% ─────────────────────────────────────────────────────────────────────────────

\subsection{Estatísticas Gerais}

A Tabela~\ref{tab:stats} apresenta as métricas consolidadas do processamento
completo do dataset.

\begin{table}[H]
\centering
\caption{Estatísticas do processamento completo (590 tiles).}
\label{tab:stats}
\begin{tabular}{ll}
\toprule
\textbf{Métrica} & \textbf{Valor} \\
\midrule
Tiles processados & 590 \\
Tiles com $\geq 1$ detecção & 564 (95,6\%) \\
Total de detecções & 4.393 \\
Detecções por tile & $7{,}45 \pm 4{,}87$ \\
Área total estimada & 23.037,14 u² \\
Área média por região & 72.933 px² \\
Área máxima por região & 3.852.048 px² \\
Cobertura HSV média & 13,862\% \\
Tempo total & 1.450,7 s \\
Tempo médio por tile & $60{,}1 \pm 5{,}4$ ms \\
Memória média por tile & 71,8 MB \\
Memória de pico por tile & 72,3 MB \\
\bottomrule
\end{tabular}
\end{table}

\subsection{Distribuição de Detecções por Tile}

\begin{figure}[H]
    \centering
    \includegraphics[width=0.8\textwidth]{../output/hist_deteccoes_por_tile.png}
    \caption{Histograma da distribuição do número de detecções por tile.}
    \label{fig:hist_det}
\end{figure}

A distribuição (Figura~\ref{fig:hist_det}) é assimétrica à direita: a maioria dos
tiles apresenta poucas detecções (1–5), enquanto um pequeno número de tiles
concentra muitas detecções, sugerindo que a \textit{Cecropia} ocorre em manchas
densas em certas regiões do mosaico.

\subsection{Distribuição das Áreas de Detecção}

\begin{figure}[H]
    \centering
    \includegraphics[width=0.8\textwidth]{../output/hist_areas_deteccao.png}
    \caption{Histograma das áreas das regiões detectadas (em px²).}
    \label{fig:hist_area}
\end{figure}

A grande maioria das detecções (Figura~\ref{fig:hist_area}) possui área próxima ao
limiar mínimo de 10.000 px², com uma cauda longa de regiões muito grandes. Regiões
extremamente grandes ($> 1 \times 10^6\;\text{px}^2$) indicam provavelmente
agrupamentos de vegetação com cor similar fundidos pelo fechamento morfológico,
e não copas individuais — um sinal de falsos positivos por agregação.

\subsection{Top 10 Tiles com Mais Detecções}

\begin{figure}[H]
    \centering
    \includegraphics[width=0.8\textwidth]{../output/top10_tiles.png}
    \caption{Top 10 tiles por número de detecções.}
    \label{fig:top10}
\end{figure}

\begin{table}[H]
\centering
\caption{Top 10 tiles por número de detecções.}
\label{tab:top10}
\begin{tabular}{lrrrr}
\toprule
\textbf{Tile} & \textbf{Detecções} & \textbf{Área (px²)} & \textbf{Área (u²)} & \textbf{Cobertura} \\
\midrule
tile\_0399.jpg & 28 & 1.005.691 & 72,31 & 26,14\% \\
tile\_0390.jpg & 23 & 1.391.146 & 100,03 & 32,44\% \\
tile\_0212.jpg & 22 & 1.208.086 & 86,86 & 30,21\% \\
tile\_0305.jpg & 21 & 1.537.900 & 110,58 & 36,17\% \\
tile\_0328.jpg & 21 & 1.129.846 & 81,24 & 28,60\% \\
tile\_0681.jpg & 21 &   454.332 & 32,67 & 14,32\% \\
tile\_0685.jpg & 21 &   714.738 & 51,39 & 21,85\% \\
tile\_0727.jpg & 21 & 1.865.179 & 134,11 & 41,95\% \\
tile\_0744.jpg & 21 &   604.596 & 43,47 & 18,47\% \\
tile\_0647.jpg & 20 &   444.228 & 31,94 & 14,49\% \\
\bottomrule
\end{tabular}
\end{table}

\subsection{Cobertura HSV por Tile}

\begin{figure}[H]
    \centering
    \includegraphics[width=0.8\textwidth]{../output/cobertura_percentual.png}
    \caption{Percentual de cobertura HSV por tile (pixels dentro da faixa HSV
             definida, após morfologia).}
    \label{fig:cobertura}
\end{figure}

\subsection{Desempenho Computacional}

\begin{figure}[H]
    \centering
    \begin{subfigure}[b]{0.48\textwidth}
        \centering
        \includegraphics[width=\textwidth]{../output/tempo_processamento.png}
        \caption{Tempo de processamento por tile.}
        \label{fig:tempo}
    \end{subfigure}
    \hfill
    \begin{subfigure}[b]{0.48\textwidth}
        \centering
        \includegraphics[width=\textwidth]{../output/memoria_por_tile.png}
        \caption{Uso de memória por tile.}
        \label{fig:memoria}
    \end{subfigure}
    \caption{Métricas de desempenho computacional do pipeline.}
    \label{fig:desempenho}
\end{figure}

O tempo médio de 60,1 ms por tile demonstra que o pipeline clássico é altamente
eficiente. Com 3 workers em paralelo (\texttt{multiprocessing.Pool}), o total de
590 tiles foi processado em $\approx$ 24 minutos. O uso de memória é estável
($\sim$72 MB/tile), reflexo das operações matriciais do OpenCV que mantêm apenas
a imagem corrente em memória.

% ─────────────────────────────────────────────────────────────────────────────
\section{Interpretação e Limitações}
% ─────────────────────────────────────────────────────────────────────────────

\subsection{Alta taxa de tiles com detecções (95,6\%)}

O fato de 564 em 590 tiles apresentarem ao menos uma detecção sugere cobertura
ampla do mosaico por vegetação dentro da faixa HSV definida. Entretanto, esse
número deve ser interpretado com cautela: o filtro HSV não é específico para
\textit{Cecropia} — outras espécies com folhagem de tonalidade verde-clara
similar também ativam a máscara.

\subsection{Falsos positivos por confusão espectral}

A inspeção visual do \texttt{tile\_0399.jpg} (Figura~\ref{fig:pipeline}) evidencia
o principal problema do método: a segmentação HSV não distingue \textit{Cecropia}
de outras plantas com tom similar. No canal H (matiz), praticamente toda a
vegetação do tile cai dentro do intervalo $[44, 58]$, fazendo a máscara bruta
ativar de forma generalizada. Após as operações morfológicas, grandes blobs são
formados fundindo múltiplas espécies, e cada blob é contado como uma detecção.

As causas raízes são:

\begin{itemize}
    \item \textbf{Sobreposição espectral:} o intervalo HSV da \textit{Cecropia}
    não é exclusivo — vegetação secundária densa tem cor similar.
    \item \textbf{Ausência de informação de forma:} a \textit{Cecropia} tem copa
    arredondada com padrão radial característico, mas o pipeline não explora
    textura nem geometria.
    \item \textbf{Fusão morfológica:} o CLOSE com kernel de 25 px une copas
    próximas (de espécies distintas) em uma única região, inflando o número de
    detecções por área.
\end{itemize}

\subsection{Regiões de área muito grande}

A área máxima de $\approx 3{,}85 \times 10^6\;\text{px}^2$ representa uma região
de $\approx 87 \times 87\;\text{m}^2$ — claramente não uma copa individual, mas
um cluster de vegetação fundido pelo fechamento morfológico. Isso indica que o
limiar de área mínima (10.000 px²) elimina ruído pequeno, mas não impede que
agregações grandes sejam contadas erroneamente como detecções únicas.

% ─────────────────────────────────────────────────────────────────────────────
\section{Conclusões}
% ─────────────────────────────────────────────────────────────────────────────

O pipeline HSV clássico demonstra ser viável para uma primeira estimativa da
distribuição espacial de \textit{Cecropia} em mosaicos aéreos: é eficiente
($\sim$60 ms/tile), escalável via paralelismo e produz resultados interpretáveis
pelas visualizações passo a passo. O total de 4.393 detecções em 590 tiles
indica densa presença de vegetação com tonalidade compatível com a espécie.

Contudo, a ausência de validação com ground truth impede afirmar quantas dessas
detecções correspondem efetivamente a indivíduos de \textit{Cecropia}. A análise
visual sugere taxa de falsos positivos relevante em tiles com vegetação densa e
heterogênea.

\subsection*{Melhorias possíveis}

Para aumentar a precisão sem recorrer a deep learning, as seguintes extensões
são recomendadas:

\begin{enumerate}
    \item \textbf{Filtro de circularidade:} calcular a razão
    $4\pi A / P^2$ (onde $A$ é a área e $P$ o perímetro do contorno) e descartar
    detecções com baixa circularidade — copas de \textit{Cecropia} tendem a ser
    mais circulares que arbustos irregulares.
    \item \textbf{Refinamento do intervalo HSV:} amostrar pixels de indivíduos
    anotados manualmente para ajustar os limiares com base em dados reais, reduzindo
    a sobreposição espectral com outras espécies.
    \item \textbf{Limiar de área máxima:} rejeitar regiões acima de um área máxima
    plausível para uma copa individual, evitando que clusters sejam contados como
    detecção única.
    \item \textbf{Análise de textura:} incorporar descritores de textura (LBP,
    Haralick) para diferenciar a estrutura radial das folhas de \textit{Cecropia}
    de outras copas.
\end{enumerate}

\end{document}
